\section{问题重述}

\subsection{问题背景}

微网是为解决光伏发电等分布式能源本地消纳、减少并网问题而发展起来的小型发用电系统。出于经济考虑，非孤岛式微网可以在外网电价较低或分布式能源有剩余的时间段为储能设备充电，并在外网电价较高的时间段释放电能。调控的目标是利用小区负载、光伏发电量与储能当前储电量等数据，尽可能精准地给出微网购电量与储能充放电量，使微网既满足小区负载，又尽可能节省购电费用。

\subsection{储能设备参数}

储能最大容量 $12000\,\mathrm{kWh}$，允许储电量范围 $1200\sim10800\,\mathrm{kWh}$，最大充放电功率 $5000\,\mathrm{kW}$，充放电效率 $90\%$（充入 $1\,\mathrm{kWh}$ 有 $0.9\,\mathrm{kWh}$ 进入电池，放出 $1\,\mathrm{kWh}$ 电池减少 $1\,\mathrm{kWh}$）。2025 年 1 月 1 日 0:00 的电量为 $6000\,\mathrm{kWh}$。调度间隔为 10 分钟，每天 144 个区间。

\subsection{需要解决的问题}

\textbf{问题一}：每天的电价和小区负载相同，并给出一天的光伏发电功率预测，要求微网提供的电能不低于小区负载，且储能设备在 0:00 与 24:00 的储电量相同。建立数学模型，给出指定 10 分钟区间的购电量、全天购电量与购电费、六个 4 小时时段的充放电量及首尾储电量。

\textbf{问题二}：每天电价相同，但小区负载和光伏发电功率随日期和时间变化。若正常购电、光伏和储能仍不能满足负载，需按交易时刻电价的 5 倍紧急购电；除此以外其他购电费用按计划购电量计算。每天 0:00 制定当天计划，对 2025 年 2 月 1 日至 12 月 31 日给出计划购电量、充放电量、连续紧急购电区间及电量。

\textbf{问题三}：每天 0:00、6:00、12:00、18:00 可获得未来 24 小时整点的光伏发电功率预报。0:00 的预报用于制定初始计划，后续预报可用于调整尚未执行的购电计划。计划购电量高于调整购电量的部分按交易时刻电价的 $50\%$ 承担违约费用；调整购电量高于计划购电量的部分按交易时刻电价的 $1.5$ 倍购入。总费用包括计划购电费用、紧急购电费用和调整相关费用。要求给出完整的初始计划、调整计划、最终储能调度与紧急购电，并分析是否需要引入其他时刻的预报。

\textbf{问题四}：外网电价按日期和 10 分钟区间实时波动，使用附件 4 的数据重新计算问题二与问题三，并给出相应成本与策略的比较。
